% 12_body.tex — Day3 산출물 2/5 (⑫) 성과보고 — EVM · Earned Schedule (빈 템플릿 / 완성 예시 공용 본문)
% 드라이버: 12_성과보고_빈템플릿.tex (\wbblanktrue) / 12_성과보고_완성예시.tex (\wbblankfalse)
% 내용 출처: PM트랙/Day3_워크북.md §2.3(항목 가이드) §2.4(빈 템플릿) §2.5(온담 P1 완성 예시본)
\documentclass[10pt,a4paper]{article}
\usepackage[a4paper,margin=16mm,top=14mm,bottom=20mm]{geometry}
\def\DocNum{2}
\def\DocTitle{성과보고 — EVM · Earned Schedule}
\def\DocSub{⑫ Performance Report (EVM · ES · Highlight Report)}
\def\DocVariant{\B{빈 템플릿}{온담 P1 완성 예시}}
\usepackage{wbstyle}
\begin{document}

\docheader
 {\Bg{[정식 명칭과 코드]}{차세대 주문·재고 통합 플랫폼 구축 (ONE ONDAM P1)}}
 {\Bg{[제n호 / 데이터 일자 / 발행일]}{제9호 / 데이터 2026-09-30 / 발행 2026-10-05}}
 {\Bg{[P1 PM · PMO 기준선 담당]}{P1 PM, P1 PMO 기준선 담당\rd{7mm}}}
 {\Bg{[스폰서 실명]}{정미란 CFO (스폰서)\rd{7mm}}}

% ------------------------------------------------------------------ 1
\secbar{1. 문서 정보와 측정 규칙}
\guide{데이터 일자와 발행일을 구분한다. BAC는 \textbf{시간대별로 배분된 PMB의 합}이고, 예외 판정 기준은 집행 한도(BAC + 컨틴전시) — 두 숫자의 관계를 명시한다. AC 인식 기준은 계약별로 다르므로 항목별로 적는다(SI 인수 대금 / 라이선스 인보이스 / 용역 월 청구 / 클라우드 사용량 / PMO 인건비). 승인되었으나 기준선에 미반영된 변경(PV 없는 AC)은 별도 행.}
{\footnotesize
\begin{tabularx}{\linewidth}{|L{30mm}|Y|L{22mm}|L{48mm}|}\hline
\hd{항목} & \hd{내용} & \hd{항목} & \hd{내용} \\ \hline
\B{\rd{10mm}}{}데이터 일자 / 발행일 & \Bg{[YYYY-MM-DD / YYYY-MM-DD — 제n호]}{2026-09-30(수) / 2026-10-05(월) — 월 성과보고 제9호} & 비교 기준선 & \Bg{[BL-0n; 다음 BL 발행 예정]}{\textbf{BL-01}(G1 2026-04-30). BL-02(CCB 10-15 결과)는 10-23 발행 예정 → 이번 호는 BL-01 기준, 승인·미반영 변경은 별도 행} \\ \hline
\B{\rd{10mm}}{}BAC / 예외 판정 기준 & \Bg{[BAC = PMB 합 / 집행 한도 = BAC + 컨틴전시 — 초과 전망 시 예외 보고]}{\textbf{BAC = PMB 145.80}(PV 곡선 합계, 컨틴전시 제외) / \textbf{집행 한도 156.00 = BAC + 컨틴전시 10.20} — EAC가 156.00 초과 전망이면 원가 예외 보고(헌장 v1.1 §12, CR-10)} & 단위 & \Bg{[억 원 소수 둘째 / 지수 셋째 / 월·영업일]}{억 원, 소수 둘째 자리; 지수 소수 셋째 자리; 기간 월·영업일} \\ \hline
\B{\rd{14mm}}{}AC 인식 기준 (발생주의) & \Bg{[항목별: SI / 상용SW / 인프라 / 전환 / 시험·감리 / PMO]}{SI = 스프린트 리뷰 인수 FP × 55만 원 + 승인된 변경 대가·대기 비용 발생분 / 상용SW = 인보이스 발생액(2026-03-13 계약 변경 후 인도 50\%·통합시험 통과 50\%) / 인프라 = 선약정 계약 시 100\%, 클라우드 월 사용료, 장비 설치 확인 / 전환 = 용역 월 청구(고정가) / 시험·감리 = 회차 보고서 접수 / PMO = 월 인건비·도구} & 미반영 승인 변경 & \Bg{[CR-nn 금액 — PV 없는 AC]}{CR-03 0.12(9월 AC 발생, PV 없음), CR-07 0.22(승인·미집행), CR-06 0.32(승인·미집행 — 재산정 공수 0.15는 6월 AC)} \\ \hline
\B{\rd{10mm}}{}작성 / 검토 / 수신 & \Bg{[작성 / 대조자 / 수신자]}{P1 PM, P1 PMO 기준선 담당 / 이재현(SI 실적 대조), 재경팀장(AC 대조) / 정미란 스폰서, 프로그램 운영위(10-08), CCB(10-15), 윤태호 사본} & 검산 스크립트 & \Bg{[경로]}{../그림/PM\_Day3/src/day3\_evm.py} \\ \hline
\end{tabularx}}

% ------------------------------------------------------------------ 2 (가로)
\landscapeon
\secbar{2. Rules of Credit — EV 측정 규칙}
\guide{WBS L2·원가 항목별로 EV를 어떻게 인정하는지 — 기법(0/100, 50/50, 마일스톤 가중, 물리적 \%, 배분 노력, LOE), 단위, 인정 시점·증빙, 진행 중 항목의 처리. Day2 ⑨의 PV 배분 규칙과 나란히 놓고 \textbf{시차}를 명시한다(PV는 스프린트 시작 월, EV는 인수 시점 → 데이터 일자에 진행 중인 스프린트 1개 분의 구조적 SV 음수). 규칙은 G1에 확정하고 실행 중 바꾸지 않는다. 담당자가 보고하는 “진척률 \%”는 규칙이 아니다.}
{\footnotesize
\begin{tabularx}{\linewidth}{|L{60mm}|L{22mm}|L{38mm}|L{42mm}|L{26mm}|Y|}\hline
\hd{WBS L2 / 원가 항목} & \hd{기법} & \hd{단위} & \hd{인정 시점 · 증빙} & \hd{진행 중 처리} & \hd{PV 규칙(Day2 ⑨)과의 시차} \\ \hline
\ifwbblank
\rd{8mm}\g{[SI — 설계 스프린트]} & \g{[0/100]} & \g{[스프린트 (억/회)]} & \g{[리뷰 인수 서명]} & \g{[0]} & \g{[PV 시작 월 / EV 인수 월]} \\ \hline
\rd{8mm}\g{[SI — 구현 스프린트]} & \g{[물리적 \%]} & \g{[인수 FP ÷ 계획 FP × 억]} & & \g{[진행 중 0]} & \\ \hline
\rd{8mm}\g{[SI — 통합·시험]} & \g{[마일스톤 가중]} & & & & \\ \hline
\rd{8mm}\g{[상용SW — 라이선스]} & \g{[마일스톤]} & \g{[인도 확인서 100\%]} & & & \g{[AC 지급 조건과의 차이]} \\ \hline
\rd{8mm}\g{[상용SW — 유지보수]} & \g{[LOE]} & & & & \\ \hline
\rd{8mm}\g{[인프라 — 선약정 / 사용료 / 장비]} & & & & & \\ \hline
\rd{8mm}\g{[단말]} & \g{[가중: 입고 / 설치]} & & & & \\ \hline
\rd{8mm}\g{[데이터 전환 — 매핑 / 리허설]} & & & & & \\ \hline
\rd{8mm}\g{[시험·감리]} & \g{[0/100 회차]} & & & & \g{[선행 실시 시 SV 양수]} \\ \hline
\rd{8mm}\g{[PMO]} & \g{[LOE]} & & & & \\ \hline
\else
SI — 설계 스프린트 S1–S6 (1.1.1·1.2.1·1.3.1·1.4.1·1.5.1) & 0/100 & 스프린트 (1.25/회) & 스프린트 리뷰에서 발주사 인수 서명 & 0 & PV는 시작 월, EV는 인수 월 → 월말에 걸친 스프린트 1개 분 시차 \\ \hline
SI — 구현 스프린트 S7–S26 (1.1.2~1.1.4, 1.2.2~1.2.5, 1.3.2~1.3.5, 1.4.2~1.4.4, 1.5.2~1.5.5) & 물리적 \% (FP) & 인수 FP ÷ 계획 620 × 1.96 & 스프린트 리뷰 인수 FP(FP 확인 3인: PMO·이재현·형상) & \textbf{진행 중 스프린트 0} & 동일 — 데이터 일자에 진행 중인 스프린트의 PV(1.96)는 구조적 SV 음수 \\ \hline
SI — 통합·시험 스프린트 S27–S36 & 마일스톤 가중 & 통합시험 TC 통과율(1.6.4) 60\% / UAT 종료 판정(1.6.7) 40\% (2.14/회 환산) & 시험 결과 보고서 & 통과율 비례 & — (2026-12 이후) \\ \hline
상용SW — 라이선스 (1.5.2·1.5.3·1.5.4) & 마일스톤 & 인도·설치 확인서 100\% & 인도 확인서(박세연) & 0 & PV는 계획 인도 월 — 인도 지연 시 SV 음수, \textbf{AC는 50/50 지급 조건으로 EV보다 늦게 발생} \\ \hline
상용SW — 1년차 유지보수 & LOE & 월 0.2 & 계약 기간 & — & 없음 \\ \hline
인프라 — 클라우드 선약정 & 마일스톤 & 계약 시 100\% (6.0) & 계약서(강현도) & 0 & 없음 \\ \hline
인프라 — 클라우드 사용료 & LOE & 월 0.4 & 사용량 청구 & — & AC는 실사용(0.5) \\ \hline
인프라 — 망분리·보안·DR 장비 & 마일스톤 & 설치 확인 100\% & 설치 확인서 & 0 & 없음 \\ \hline
단말 (1.1.5) & 마일스톤 가중 & 입고 검수 70\% / 설치 확인 30\% & 검수 조서·설치 확인서 & 0 & 없음 (2026-12 이후) \\ \hline
데이터 전환 — 매핑·정제 (1.5.7·1.5.8) & 물리적 \% & 매핑 확정 테이블 수 ÷ 1,842 × 구간 가치 & 매핑 정의서 승인(박세연) & 비례 & 없음 \\ \hline
데이터 전환 — 리허설 (1.5.9) & 0/100 & 회차 (3회) & 검증 6종 보고서 & 0 & 없음 \\ \hline
시험·감리 (1.6.4~1.6.8) & 0/100 & 회차 보고서 (감리 G2 1.4 = 중간 1차 0.7 + 2차 0.7) & 보고서 접수 & 0 & PV는 게이트 월 — 선행 실시 시 SV 양수(진척 아님) \\ \hline
PMO (1.6.2·1.6.3) & LOE & 월 0.56 & 월 & — & 없음 \\ \hline
\fi
\end{tabularx}}
\B{}{\small G1 확정, 원가 기준선 v1.0 부속 2 [추가 설정]. LOE 항목(PMO·클라우드 사용료·유지보수)은 EV ≡ PV — 이 항목의 CV만 의미가 있고 SV는 항상 0이다.\par}

% ------------------------------------------------------------------ 3 (가로)
\clearpage
\secbar{3. 월별 PV · EV · AC}
\guide{PV는 기준선(Day2 ⑨ PV 표) 그대로 — 실적에 맞춰 고치지 않는다(재기준선 없이). EV·AC 월 값은 아래 항목별 표의 합이며 소수 둘째 자리 반올림으로 ±0.01 차이 가능. 누적만 있고 월별이 없으면 언제 무엇이 일어났는지 모른다. EV와 AC를 같은 값으로 적지 않는다(고정가 착시를 전 항목에 적용하는 오류).}
{\footnotesize\setlength{\tabcolsep}{3pt}
\begin{tabularx}{\linewidth}{|L{22mm}|C{16mm}|C{16mm}|C{16mm}|C{18mm}|C{18mm}|C{18mm}|C{16mm}|C{16mm}|C{16mm}|C{16mm}|Y|}\hline
\hd{월} & \hd{PV} & \hd{EV} & \hd{AC} & \hd{PV 누적} & \hd{EV 누적} & \hd{AC 누적} & \hd{SV} & \hd{CV} & \hd{SPI} & \hd{CPI} & \hd{비고} \\ \hline
\ifwbblank
\rd{6mm}\g{[YYYY-MM]} & & & & & & & \g{[EV−PV]} & \g{[EV−AC]} & \g{[EV/PV]} & \g{[EV/AC]} & \\ \hline
\rd{6mm} & & & & & & & & & & & \\ \hline
\rd{6mm} & & & & & & & & & & & \\ \hline
\rd{6mm} & & & & & & & & & & & \\ \hline
\rd{6mm} & & & & & & & & & & & \\ \hline
\rd{6mm} & & & & & & & & & & & \\ \hline
\rd{6mm} & & & & & & & & & & & \\ \hline
\rd{6mm} & & & & & & & & & & & \\ \hline
\rd{6mm} & & & & & & & & & & & \\ \hline
\rd{6mm} & & & & & & & & & & & \\ \hline
\rd{6mm} & & & & & & & & & & & \\ \hline
\rd{6mm} & & & & & & & & & & & \\ \hline
\rd{6mm}\textbf{데이터 일자} & & & & & & & & & & & \\ \hline
\else
2026-01 & 3.06 & 1.81 & 1.85 & 3.06 & 1.81 & 1.85 & −1.25 & −0.04 & 0.592 & 0.978 & S1 인수 \\ \hline
2026-02 & 3.06 & 3.06 & 3.10 & 6.12 & 4.87 & 4.95 & −1.25 & −0.08 & 0.796 & 0.984 & S2·S3 인수 \\ \hline
2026-03 & 13.02 & 12.31 & 12.35 & 19.14 & 17.18 & 17.30 & −1.96 & −0.12 & 0.898 & 0.993 & API GW 8.0 인도 \\ \hline
2026-04 & 17.69 & 11.68 & 11.72 & 36.83 & 28.86 & 29.02 & −7.97 & −0.16 & 0.784 & 0.994 & 상용SW −6.0(조달 이연) \\ \hline
2026-05 & 8.89 & 10.97 & 8.11 & 45.72 & 39.83 & 37.13 & −5.89 & +2.70 & 0.871 & 1.073 & 라이선스 인도, AC 50\% \\ \hline
2026-06 & 8.89 & 8.27 & 8.81 & 54.61 & 48.10 & 45.94 & −6.51 & +2.16 & 0.881 & 1.047 & R-05 실현 여파 \\ \hline
2026-07 & 8.09 & 9.18 & 7.99 & 62.69 & 57.28 & 53.93 & −5.41 & +3.35 & 0.914 & 1.062 & SAP 애드온 인도 \\ \hline
2026-08 & 9.05 & 8.23 & 9.02 & 71.74 & 65.51 & 62.95 & −6.23 & +2.56 & 0.913 & 1.041 & 규제 검증 집중 \\ \hline
\textbf{2026-09} & 6.09 & 7.92 & 8.18 & \textbf{77.83} & \textbf{73.43} & \textbf{71.13} & \textbf{−4.40} & \textbf{+2.30} & \textbf{0.943} & \textbf{1.032} & 데이터 일자 \\ \hline
\fi
\end{tabularx}}
\B{}{\small PV 월 값·누적은 Day2 §4.5 표 그대로. 월별 SPI 추이(누적): 0.592 → 0.796 → 0.898 → \textbf{0.784(4월)} → 0.871 → 0.881 → 0.914 → 0.913 → 0.943. 4월 하락의 분해: 상용SW −6.00(MDM 2차·APM 인도 5월로 이연 — R-04 대응), SI −1.97(S9 05-10 인수 시차 — 실질 0), 인프라 0. \textbf{4월 SPI는 조달 지연이고 개발 지연이 아니다.}\par}

\needspace{55mm}\vspace{2mm}\noindent\textbf{\small 항목별 월 EV} \g{(억 원 — 누적 EV·누적 PV·SV를 항목별로 분해. SI EV의 근거는 인수 FP)}\par\vspace{0.5mm}
{\footnotesize\setlength{\tabcolsep}{3pt}
\begin{tabularx}{\linewidth}{|L{26mm}|C{13mm}|C{13mm}|C{13mm}|C{13mm}|C{13mm}|C{13mm}|C{13mm}|C{13mm}|C{13mm}|C{17mm}|C{17mm}|C{16mm}|Y|}\hline
\hd{항목} & \hd{01} & \hd{02} & \hd{03} & \hd{04} & \hd{05} & \hd{06} & \hd{07} & \hd{08} & \hd{09} & \hd{누적 EV} & \hd{누적 PV} & \hd{SV} & \hd{원인 한 줄} \\ \hline
\ifwbblank
\rd{6mm}SI & & & & & & & & & & & & & \\ \hline
\rd{6mm}상용SW & & & & & & & & & & & & & \\ \hline
\rd{6mm}인프라 & & & & & & & & & & & & & \\ \hline
\rd{6mm}단말 & & & & & & & & & & & & & \\ \hline
\rd{6mm}데이터 전환 & & & & & & & & & & & & & \\ \hline
\rd{6mm}시험·감리 & & & & & & & & & & & & & \\ \hline
\rd{6mm}PMO & & & & & & & & & & & & & \\ \hline
\rd{6mm}\textbf{합계} & & & & & & & & & & & & & \\ \hline
\else
SI & 1.25 & 2.50 & 3.75 & 3.92 & 3.41 & 3.26 & 3.67 & 3.92 & 4.01 & \textbf{29.69} & 35.00 & \textbf{−5.31} & S10 3주(R-05), R2 이월 420 FP, S20 시차 1.96 \\ \hline
상용SW & — & — & 8.00 & — & 6.00 & 0.20 & 3.60 & 0.20 & 0.20 & 18.20 & 18.20 & 0.00 & 4월 −6.0 → 5·7월 인도로 회복 \\ \hline
인프라 & — & — & — & 6.40 & 0.40 & 3.40 & 0.40 & 2.50 & 0.40 & 13.50 & 13.50 & 0.00 & 선약정 04-24 \\ \hline
단말 & — & — & — & — & — & — & — & — & — & 0.00 & 0.00 & 0.00 & 12월 이후 \\ \hline
데이터 전환 & — & — & — & — & 0.60 & 0.85 & 0.95 & 1.05 & 1.05 & 4.50 & 4.30 & +0.20 & 매핑 진도 선행(O-01) \\ \hline
시험·감리 & — & — & — & 0.80 & — & — & — & — & 1.70 & 2.50 & 1.80 & +0.70 & 중간감리 1차 선행 실시(시점 이동) \\ \hline
PMO & 0.56 & 0.56 & 0.56 & 0.56 & 0.56 & 0.56 & 0.56 & 0.56 & 0.56 & 5.04 & 5.04 & 0.00 & LOE \\ \hline
\textbf{합계} & 1.81 & 3.06 & 12.31 & 11.68 & 10.97 & 8.27 & 9.18 & 8.23 & 7.92 & \textbf{73.43} & \textbf{77.83} & \textbf{−4.40} & \\ \hline
\fi
\end{tabularx}}
\B{\small\g{SI EV 근거 — 월별 인수 FP: \hrulefill}\par}{\small SI EV 근거 — 인수 FP(스프린트 리뷰 인수 월 기준): 4월 1,240(S7 620·S8 620) / 5월 1,080(S9 600·S10 480) / 6월 1,030(S11 470·S12 560) / 7월 1,160(S13 560·S14 600) / 8월 1,240(S15 620·S16 620) / 9월 1,270(S17 640·S18 630) = \textbf{누적 7,020 FP(56.6\%)}. 계획(S7–S18 12회 × 620) 7,440 → 이월 420. S19(09-28 착수) 진행 중 0. 설계 스프린트 인수: 1월 S1 / 2월 S2·S3 / 3월 S4·S5·S6.\par}

\needspace{55mm}\vspace{2mm}\noindent\textbf{\small 항목별 월 AC} \g{(억 원 — 발생주의. CV의 원인을 항목별로 한 줄)}\par\vspace{0.5mm}
{\footnotesize\setlength{\tabcolsep}{3pt}
\begin{tabularx}{\linewidth}{|L{26mm}|C{13mm}|C{13mm}|C{13mm}|C{13mm}|C{13mm}|C{13mm}|C{13mm}|C{13mm}|C{13mm}|C{17mm}|C{17mm}|C{16mm}|Y|}\hline
\hd{항목} & \hd{01} & \hd{02} & \hd{03} & \hd{04} & \hd{05} & \hd{06} & \hd{07} & \hd{08} & \hd{09} & \hd{누적 AC} & \hd{누적 EV} & \hd{CV} & \hd{원인} \\ \hline
\ifwbblank
\rd{6mm}SI & & & & & & & & & & & & & \\ \hline
\rd{6mm}상용SW & & & & & & & & & & & & & \\ \hline
\rd{6mm}인프라 & & & & & & & & & & & & & \\ \hline
\rd{6mm}단말 & & & & & & & & & & & & & \\ \hline
\rd{6mm}데이터 전환 & & & & & & & & & & & & & \\ \hline
\rd{6mm}시험·감리 & & & & & & & & & & & & & \\ \hline
\rd{6mm}PMO & & & & & & & & & & & & & \\ \hline
\rd{6mm}\textbf{합계} & & & & & & & & & & & & & \\ \hline
\else
SI & 1.25 & 2.50 & 3.75 & 3.92 & 3.41 & 3.41 & 4.09 & 4.42 & 4.13 & \textbf{30.88} & 29.69 & \textbf{−1.19} & 인수 대금 = EV + 컨틴전시 집행 4건(6월 R-07 재산정 0.15, 7월 R-05 대기 비용 0.42, 8월 R-03 폴백 재시험 0.50, 9월 CR-03 0.12) \\ \hline
상용SW & — & — & 8.00 & — & 3.00 & 0.20 & 1.90 & 0.20 & 0.20 & 13.50 & 18.20 & \textbf{+4.70} & MDM 2차·APM 6.0 인도(5월) 50\%, SAP 애드온 3.4 인도(7월) 50\% — \textbf{미발생 4.70은 통합시험 통과(2027-01) 시 발생} \\ \hline
인프라 & — & — & — & 6.40 & 0.50 & 3.80 & 0.50 & 2.80 & 0.50 & 14.50 & 13.50 & −1.00 & 클라우드 개발환경 월 0.5(계획 0.4, 6개월 +0.6), 망분리 3.3(+0.3)·2.3(+0.1) \\ \hline
단말 & — & — & — & — & — & — & — & — & — & 0.00 & 0.00 & 0.00 & — \\ \hline
데이터 전환 & — & — & — & — & 0.60 & 0.80 & 0.90 & 1.00 & 1.00 & 4.30 & 4.50 & +0.20 & 고정가 월 청구 = 계획, 진도 선행 \\ \hline
시험·감리 & — & — & — & 0.80 & — & — & — & — & 1.75 & 2.55 & 2.50 & −0.05 & 보안진단 1회차 1.05(+0.05), 중간감리 1차 0.70 \\ \hline
PMO & 0.60 & 0.60 & 0.60 & 0.60 & 0.60 & 0.60 & 0.60 & 0.60 & 0.60 & 5.40 & 5.04 & −0.36 & 월 0.60(계획 0.56 — 도구 라이선스 증분·회의 증가) \\ \hline
\textbf{합계} & 1.85 & 3.10 & 12.35 & 11.72 & 8.11 & 8.81 & 7.99 & 9.02 & 8.18 & \textbf{71.13} & \textbf{73.43} & \textbf{+2.30} & \\ \hline
\fi
\end{tabularx}}
\B{\small\g{미반영 승인 변경 (PV 없는 AC): \hrulefill}\par}{}

% ------------------------------------------------------------------ 4 (가로)
\clearpage
\secbar{4. 편차와 지수 — SV · CV · SPI · CPI}
\guide{SV = EV − PV, CV = EV − AC, SPI = EV/PV, CPI = EV/AC. 조정값은 별도 행(조정 사유 명시)으로 \textbf{병기}하며 본값을 대체하지 않는다. 항목별 분해(SI 열만·상용SW 열만)를 붙인다. “SPI 0.943 = 약 6\% 지연”으로 읽지 않는다 — SPI는 비율이지 일수가 아니다(항목 6).}
{\footnotesize
\begin{tabularx}{\linewidth}{|L{44mm}|C{18mm}|C{18mm}|C{18mm}|C{18mm}|Y|}\hline
\hd{구분} & \hd{SV} & \hd{CV} & \hd{SPI} & \hd{CPI} & \hd{조정 사유} \\ \hline
\ifwbblank
\rd{8mm}\textbf{본값 (BL-0n 규칙)} & & & & & — \\ \hline
\rd{8mm}조정 1: \g{[측정 시차 제거]} & & & & & \g{[진행 중 스프린트 PV 제거 근거]} \\ \hline
\rd{8mm}조정 2: \g{[인식 기준 정합]} & & & & & \g{[미발생 AC 가산 근거]} \\ \hline
\rd{8mm}조정 1+2 & & & & & \\ \hline
\rd{8mm}\g{[항목] 열만} & & & & & \\ \hline
\rd{8mm}\g{[항목] 열만} & & & & & \\ \hline
\else
\textbf{본값 (BL-01 규칙)} & \textbf{−4.40} & \textbf{+2.30} & \textbf{0.943} & \textbf{1.032} & — \\ \hline
조정 1: 측정 시차 제거 & −2.44 & +2.30 & 0.968 & 1.032 & S20(09-28 착수) PV 1.96은 정상 진행에서도 EV 0 — PV 77.83 − 1.96 = 75.87 기준 \\ \hline
조정 2: 발생주의 정합 & −4.40 & \textbf{−2.40} & 0.943 & \textbf{0.968} & 라이선스 미발생 4.70을 AC에 가산(AC 75.83) — 인도된 가치의 원가는 발생했다 \\ \hline
조정 1+2 & −2.44 & −2.40 & 0.968 & 0.968 & 실질 지연 약 1.7 스프린트 분, 실질 원가 초과 2.40 \\ \hline
SI 열만 & −5.31 & −1.19 & 0.848 & 0.961 & 고정가 — CV 전액이 컨틴전시 집행 \\ \hline
상용SW 열만 & 0.00 & +4.70 & 1.000 & 1.348 & 인식 기준 차이 \\ \hline
\fi
\end{tabularx}}
\B{\small\g{월별 SPI 추이와 최저점의 항목별 분해: \hrulefill}\par}{\small 5월 이후 SI SV: −2.49 → −3.16 → −3.42 → −5.39 → −5.31 — R-05 실현(S10 3주)과 R2 규제 항목 전진 배치의 이월이 누적. 스폰서가 읽어야 할 한 줄: 조정 1+2 후 SPI 0.968·CPI 0.968 — “일정도 원가도 약 3\% 뒤”. 그러나 3\%는 평균이고 지연은 임계경로에 몰려 있다(항목 8).\par}

% ------------------------------------------------------------------ 5 (가로)
\secbar[70mm]{5. 예측 — EAC 4공식 · TCPI · VAC}
\guide{EAC 네 가지와 각각의 전제, 어느 것을 공식 EAC로 채택했는지와 이유. 상향식 ETC는 항목별 근거(잔여 계약·미발생·변경 대가·견적). 조정 CPI로 계산한 값을 병기한다. VAC = BAC − EAC는 승인 예산(168.0)이 아니라 BAC·집행 한도와 비교한다. TCPI는 BAC 기준과 집행 한도 기준을 구분한다.}
{\footnotesize
\begin{tabularx}{\linewidth}{|L{66mm}|L{40mm}|C{16mm}|C{20mm}|C{20mm}|Y|}\hline
\hd{EAC} & \hd{전제} & \hd{값} & \hd{VAC = BAC − EAC} & \hd{집행 한도 여유} & \hd{채택} \\ \hline
\ifwbblank
\rd{8mm}EAC1 = BAC ÷ CPI & \g{[현재 원가 효율 유지]} & & & & \\ \hline
\rd{8mm}EAC2 = AC + (BAC − EV) & \g{[잔여는 계획대로]} & & & & \\ \hline
\rd{8mm}EAC3 = AC + (BAC − EV) ÷ (CPI × SPI) & \g{[원가·일정 효율 복합]} & & & & \\ \hline
\rd{8mm}EAC4 = AC + 상향식 ETC & \g{[항목별 근거 — 아래 표]} & & & & \\ \hline
\rd{8mm}EAC1′ = BAC ÷ 조정 CPI & & & & & \\ \hline
\rd{8mm}EAC3′ = 조정 AC + (BAC − EV) ÷ (조정 CPI × SPI) & & & & & \\ \hline
\else
EAC1 = BAC ÷ CPI = 145.80 ÷ 1.032 & 현재 원가 효율이 끝까지 유지 & 141.23 & +4.57 & 14.77 & 기각 — 착시 포함, PMB보다 낮은 EAC는 고정가 구조에서 불가능 \\ \hline
EAC2 = AC + (BAC − EV) = 71.13 + 72.37 & 잔여는 계획대로 & 143.50 & +2.30 & 12.50 & 기각 — 라이선스 미발생 4.70·컨틴전시 잔여 집행 무시 \\ \hline
EAC3 = AC + (BAC − EV) ÷ (CPI × SPI) = 71.13 + 72.37 ÷ 0.973 & 원가·일정 효율 복합 & 145.43 & +0.37 & 10.57 & 참고 \\ \hline
\textbf{EAC4 = AC + 상향식 ETC = 71.13 + 79.72} & 항목별 잔여 근거(아래) & \textbf{150.86} & \textbf{−5.06} & \textbf{5.14} & \textbf{공식 EAC} \\ \hline
EAC1′ = BAC ÷ CPI(조정 0.968) & 발생주의 CPI 유지 & 150.57 & −4.77 & 5.43 & EAC4와 정합 — 검산 \\ \hline
EAC3′ = AC(조정 75.83) + 72.37 ÷ (0.968 × 0.943) & 조정 CPI × SPI & \textbf{155.04} & −9.24 & \textbf{0.96} & \textbf{경고 대역} — 일정 회복 실패 시 도착지 \\ \hline
\fi
\end{tabularx}}
\vspace{2mm}
\noindent\begin{minipage}[t]{0.5\linewidth}\vspace{0pt}
\noindent\textbf{\small 상향식 ETC의 근거}\par\vspace{0.5mm}
{\footnotesize
\begin{tabularx}{\linewidth}[t]{|L{40mm}|C{14mm}|Y|}\hline
\hd{항목} & \hd{ETC} & \hd{근거} \\ \hline
\ifwbblank
\rd{7mm}SI 잔여 계약 & & \g{[계약 총액 − EV]} \\ \hline
\rd{7mm}SI 승인·예정 변경 대가 & & \g{[FP Δ × 단가]} \\ \hline
\rd{7mm}SI 대기 비용 잔여 & & \\ \hline
\rd{7mm}상용SW & & \g{[잔여 PV + 미발생]} \\ \hline
\rd{7mm}인프라 & & \g{[잔여 PV × 추세]} \\ \hline
\rd{7mm}단말 & & \g{[견적]} \\ \hline
\rd{7mm}데이터 전환 & & \\ \hline
\rd{7mm}시험·감리 & & \\ \hline
\rd{7mm}PMO & & \g{[잔여 월 × 실적 단가]} \\ \hline
\rd{7mm}\textbf{합계} & & \\ \hline
\else
SI 잔여 계약 & 38.51 & 68.20 − EV 29.69 (잔여 FP 5,380 × 55만 원 = 29.59 + 통합·시험 스프린트 가치 — 계약 총액 기준) \\ \hline
SI 승인·예정 변경 대가 & 1.02 & +185 FP(⑪ 변경 로그, BL-02 편입 예정) \\ \hline
SI 대기 비용 잔여 추정 & 0.18 & R-05 배정 잔여(재발 대비) \\ \hline
상용SW & 5.90 & 잔여 PV 1.20(유지보수) + \textbf{미발생 4.70}(통합시험 통과 시) \\ \hline
인프라 & 7.92 & 잔여 PV 7.40 × 1.07(클라우드 실사용 추세) \\ \hline
단말 & 8.90 & 8.60 + 단일 공급사 견적 상승 0.30(강현도 09-18 [추가 설정]) \\ \hline
데이터 전환 & 5.10 & 9.40 − 4.50 + 아카이브 정책 수립 0.20 \\ \hline
시험·감리 & 7.40 & 9.70 − 2.50 + 부하시험 재견적 0.20(피크 5,100 케이스/h 시나리오 추가) \\ \hline
PMO & 4.80 & 8개월 × 0.60(실적 단가) \\ \hline
\textbf{합계} & \textbf{79.72} & \\ \hline
\fi
\end{tabularx}}
\end{minipage}\hfill
\begin{minipage}[t]{0.48\linewidth}\vspace{0pt}
\noindent\textbf{\small TCPI}\par\vspace{0.5mm}
{\footnotesize
\begin{tabularx}{\linewidth}[t]{|L{30mm}|L{44mm}|C{12mm}|Y|}\hline
\hd{TCPI} & \hd{산식} & \hd{값} & \hd{해석} \\ \hline
\ifwbblank
\rd{9mm}TCPI(BAC) & \g{(BAC − EV) ÷ (BAC − AC)} & & \\ \hline
\rd{9mm}TCPI(BAC, 조정 AC) & & & \\ \hline
\rd{9mm}TCPI(집행 한도) & \g{(BAC − EV) ÷ (집행 한도 − AC)} & & \\ \hline
\rd{9mm}TCPI(EAC) & \g{(BAC − EV) ÷ (EAC − AC)} & & \\ \hline
\else
TCPI(BAC) & (145.80 − 73.43) ÷ (145.80 − 71.13) = 72.37 ÷ 74.67 & 0.969 & 본값 기준 — 착시 포함, 의미 없음 \\ \hline
TCPI(BAC, 조정 AC) & 72.37 ÷ (145.80 − 75.83) = 72.37 ÷ 69.97 & \textbf{1.034} & PMB 안에 끝내려면 잔여 구간에서 3.4\% 절감 필요 — 고정가 구조에서 비현실적 → PMB 초과는 확정 \\ \hline
TCPI(집행 한도 156.00) & 72.37 ÷ (156.00 − 71.13) = 72.37 ÷ 84.87 & 0.853 & 집행 한도 안에 끝내는 것은 가능(컨틴전시 잔여 8.79가 완충) \\ \hline
TCPI(EAC4 150.86) & 72.37 ÷ (150.86 − 71.13) = 72.37 ÷ 79.73 & 0.908 & 공식 EAC 달성에 필요한 효율 \\ \hline
\fi
\end{tabularx}}
\vspace{2mm}\noindent\textbf{\small 원가 예외 보고 판정} \g{(불요/필요 — 근거, 경고 조건)}\par\vspace{0.5mm}
\begin{fieldboxu}\small
\ifwbblank\lines{4}{7mm}\else
\textbf{불요.} 공식 EAC4 150.86 $<$ 집행 한도 156.00(여유 5.14). 단, ① EAC3′ 155.04(여유 0.96)는 일정 회복 실패(A14·A16 압축 원가, 2차 창 재작업)가 실현될 때의 도착지이고 ② 잔여 컨틴전시 8.79 중 승인 변경 편입 예정 1.02를 빼면 7.77이며 컨틴전시 대상 잔여 EMV 5.06(⑬)과의 여유는 2.71 — \textbf{다음 보고(10-31 데이터)에서 EAC4가 153.0을 넘으면 예외 보고 준비}를 스폰서에 예고한다(경고 등급 A).
\fi
\end{fieldboxu}
\end{minipage}\par

% ------------------------------------------------------------------ 6 (가로)
\clearpage
\secbar{6. Earned Schedule — ES · SV(t) · SPI(t) · IEAC(t)}
\guide{ES는 보간으로 구한다 — PV 누적이 EV를 넘지 않는 마지막 월 n에 (EV − PV(n)) ÷ (PV(n+1) − PV(n))을 더한다(단위 월). SV(t) = ES − AT, SPI(t) = ES/AT, IEAC(t) = PD/SPI(t). IEAC(t)는 \textbf{추세 외삽}이므로 CPM 회복 계획(항목 8)과 병기하고 어느 것이 계획이고 어느 것이 추세인가를 적는다. SPI는 종료로 갈수록 1로 수렴한다 — 월별 SPI와 SPI(t)를 나란히 놓아 벌어짐을 본다.}
\noindent\begin{minipage}[t]{0.5\linewidth}\vspace{0pt}
{\footnotesize
\begin{tabularx}{\linewidth}[t]{|L{22mm}|Y|C{30mm}|}\hline
\hd{지표} & \hd{산식} & \hd{값} \\ \hline
\ifwbblank
\rd{9mm}ES & \g{n + (EV − PV(n)) ÷ (PV(n+1) − PV(n))} & \g{[개월]} \\ \hline
\rd{7mm}AT & \g{[시작월 ~ 데이터 월]} & \g{[개월]} \\ \hline
\rd{7mm}SV(t) & ES − AT & \g{[개월 ≈ 영업일]} \\ \hline
\rd{7mm}SPI(t) & ES ÷ AT & \\ \hline
\rd{7mm}PD & 계획 기간 & \g{[개월]} \\ \hline
\rd{7mm}IEAC(t) & PD ÷ SPI(t) & \g{[개월 → +n영업일]} \\ \hline
\rd{9mm}추세 종료일 & \g{[기준선 종료일 + 지연] vs 허용범위} & \\ \hline
\else
ES & 8 + (EV 73.43 − PV(8월) 71.74) ÷ (PV(9월) 77.83 − PV(8월) 71.74) = 8 + 1.69 ÷ 6.09 & \textbf{8.28개월} \\ \hline
AT & 2026-01 ~ 2026-09 & 9개월 \\ \hline
SV(t) & ES − AT & \textbf{−0.72개월} (≈ −15영업일) \\ \hline
SPI(t) & ES ÷ AT & \textbf{0.920} \\ \hline
PD & 계획 기간 & 17개월 \\ \hline
IEAC(t) & PD ÷ SPI(t) = 17 ÷ 0.920 & \textbf{18.48개월} → +1.48개월 ≈ \textbf{+31영업일} \\ \hline
추세 종료일 & 2027-05-28 + 31영업일 & \textbf{2027-07-12} [추세 외삽] vs 허용범위 2027-06-18(+15영업일) → \textbf{초과} \\ \hline
\fi
\end{tabularx}}
\end{minipage}\hfill
\begin{minipage}[t]{0.48\linewidth}\vspace{0pt}
{\footnotesize
\begin{tabularx}{\linewidth}[t]{|L{16mm}|C{13mm}|C{13mm}|C{13mm}|Y|}\hline
\hd{월} & \hd{SPI} & \hd{SPI(t)} & \hd{차이} & \hd{판독} \\ \hline
\ifwbblank
\rd{7mm}\g{[YYYY-MM]} & & & & \\ \hline
\rd{7mm} & & & & \\ \hline
\rd{7mm} & & & & \\ \hline
\rd{7mm} & & & & \\ \hline
\rd{7mm} & & & & \\ \hline
\rd{7mm} & & & & \\ \hline
\rd{7mm}\textbf{데이터 월} & & & & \g{[벌어짐의 방향]} \\ \hline
\else
2026-03 & 0.898 & 0.950 & +0.052 & 조달 봉우리 앞 — ES가 PV 완만 구간에 있어 SPI(t)가 관대 \\ \hline
2026-04 & 0.784 & 0.887 & +0.103 & 조달 이연: 돈으로는 −8.0, 시간으로는 −0.45개월 \\ \hline
2026-05 & 0.871 & 0.868 & −0.003 & 라이선스 인도로 SPI 회복, SPI(t)는 R-05 지연을 그대로 반영 \\ \hline
2026-06 & 0.881 & 0.878 & −0.003 & \\ \hline
2026-07 & 0.914 & 0.904 & −0.010 & \textbf{벌어지기 시작} \\ \hline
2026-08 & 0.913 & 0.914 & +0.001 & \\ \hline
\textbf{2026-09} & \textbf{0.943} & \textbf{0.920} & \textbf{−0.023} & SPI는 회복처럼, SPI(t)는 정체 — 수렴 착시 \\ \hline
\fi
\end{tabularx}}
\end{minipage}\par

% ------------------------------------------------------------------ 7 (가로)
\secbar[60mm]{7. 착시 해부}
\guide{지수를 좋게 또는 나쁘게 보이게 하는 구조적 원인 각각에 이름을 붙이고 — 무엇이, 왜, 얼마만큼, 조정 후 값. 조정값은 항목 4·5에 반영된 것과 같아야 한다. 조정은 병기하는 것이지 대체하는 것이 아니다. 알면서 유리한 쪽만 보고하지 않는다.}
{\footnotesize
\begin{tabularx}{\linewidth}{|C{6mm}|L{34mm}|L{40mm}|Y|L{30mm}|L{44mm}|}\hline
\hd{\#} & \hd{이름} & \hd{무엇이} & \hd{왜} & \hd{크기} & \hd{조정 후} \\ \hline
\ifwbblank
\rd{9mm}1 & \g{[인식 기준 차이]} & & & & \\ \hline
\rd{9mm}2 & \g{[고정가 CPI ≡ 1]} & & & & \\ \hline
\rd{9mm}3 & \g{[측정 시차]} & & & & \\ \hline
\rd{9mm}4 & \g{[SPI 수렴]} & & & & \\ \hline
\rd{9mm}5 & \g{[선행 인도의 양수 SV]} & & & & \\ \hline
\else
1 & \textbf{라이선스 인식 기준 차이} & 상용SW CV +4.70, 전체 CPI 1.032 & R-04 대응으로 2026-03-13 지급 조건을 인도 50\%·통합시험 통과 50\%로 변경. EV는 인도 확인서에 100\%, AC는 인보이스 50\% & +4.70 (전체 CV의 200\%) & CPI 0.968, CV −2.40 \\ \hline
2 & \textbf{고정가 SI의 CPI ≡ 1} & SI CPI 0.961 & 발주사 관점 SI AC = 인수 FP 대금 = EV. CV는 오직 컨틴전시 집행(1.19)과 변경 대가 & 구조적 & SI 열의 원가 성과는 “컨틴전시 소진율”로 읽는다(⑬ 항목 7) \\ \hline
3 & \textbf{측정 시차} & SV −4.40 중 −1.96 & PV는 스프린트 시작 월에 계상, EV는 인수 시 계상. S20(09-28 착수)의 1.96은 정상 진행에서도 EV 0 & 1.96 (SV의 45\%) & SV −2.44, SPI 0.968 — \textbf{실질 지연 = 약 1.7 스프린트} \\ \hline
4 & \textbf{SPI 수렴} & 9월 SPI 0.943 상승(8월 0.913) & 누적 EV가 커질수록 SPI는 1로 수렴하며 종료 시 정확히 1이 된다. 실제 지연은 SPI(t) 0.920(하락) & SPI − SPI(t) = +0.023, 확대 중 & 일정 보고는 SPI(t)·ES·CPM TF로 \\ \hline
5 & \textbf{선행 인도의 양수 SV} & 시험·감리 SV +0.70 & 중간감리 1차를 G2 감리(11월 PV)의 절반으로 9월에 실시. 진척이 아니라 시점 이동 & +0.70 & 진척 판단에서 제외 — 단 감리 지적 8건(⑭)은 실제 \\ \hline
\fi
\end{tabularx}}

% ------------------------------------------------------------------ 8 (가로)
\clearpage
\secbar{8. 일정 통제 — TF 소진 · 마일스톤 전망 · 버퍼 · 회복 계획}
\guide{CPM 계산표(Day2 ⑧ 항목 5)의 활동별 실제 착수·완료(또는 전망)와 TF 소진. 지연을 “n일”로만 적지 않고 TF로 본다 — 여유 20일인 활동의 5일 지연과 여유 0인 활동의 5일 지연은 다른 사건이다. 마일스톤 전망은 회복 계획 없는 값(as-is)과 있는 값(to-be)을 나란히. 컷오버 창 이탈 예상은 즉시 예외 보고. 버퍼를 P1이 소비해 놓고 프로그램에 알리지 않는 것이 흔한 오류다.}
{\scriptsize\setlength{\tabcolsep}{2.5pt}
\begin{longtable}{|L{36mm}|L{40mm}|L{78mm}|L{40mm}|L{62mm}|}\hline
\hd{활동} & \hd{기준선 ES → EF (일자)} & \hd{실제 · 전망} & \hd{TF 전 → 후} & \hd{판독} \\ \hline \endhead
\ifwbblank
\rd{6mm}A01 & & & & \\ \hline
\rd{6mm}A02 & & & & \\ \hline
\rd{6mm}A03 & & & & \\ \hline
\rd{6mm}A04 & & & & \\ \hline
\rd{6mm}A05 & & & & \\ \hline
\rd{6mm}A06 & & & & \\ \hline
\rd{6mm}A07 & & & & \\ \hline
\rd{6mm}A08 & & & & \\ \hline
\rd{6mm}A09 & & & & \\ \hline
\rd{6mm}A10 & & & & \\ \hline
\rd{6mm}A11 & & & & \\ \hline
\rd{6mm}A12 & & & & \\ \hline
\rd{6mm}A13 & & & & \\ \hline
\rd{6mm}A14 & & & & \\ \hline
\rd{6mm}A15 & & & & \\ \hline
\rd{6mm}A16 & & & & \\ \hline
\rd{6mm}MG3 & & & & \\ \hline
\rd{6mm}A17 & & & & \\ \hline
\rd{6mm}A20 & & & & \\ \hline
\else
A01 설계 S1–S6 & 0 → 60 (~03-27) & 완료 03-27 & 0 → 0 & 정시 \\ \hline
A02 13본 복원·47본 문서 & 10 → 60 & 완료 03-25 (O-01 부분 실현, 미등재 2본 발견 → CR-07) & 5 → 5 & 정시 \\ \hline
A03 조달 & 25 → 70 (~04-10) & 선약정 04-24 서명(+10), MDM 2차·APM 05-22 인도(+30), SAP 애드온 07-10(+60) & 5 → \textbf{−25(라이선스)} — 단 A05 SS+30 의존은 개발환경(선약정)만 → A05 영향 0 & R-04 대응 결과. 임계경로 미영향(A05 착수 03-30 정시) \\ \hline
A04 기준선 → G1 & 60 → 80 (~04-24) & 완료 04-24, G1 04-30 통과 & 4 → 0 & 정시 \\ \hline
A05 구현 R1 S7–S12 & 60 → 120 (~06-19) & S10 3주(05-11~05-29, R-05), R1 종료 \textbf{06-28}(+5영업일) & 0 → \textbf{−5} & 임계경로 지연 시작 \\ \hline
A06 구현 R2 S13–S18 & 120 → 180 (~09-11) & S13 06-29 착수, S18 3주(09-07~09-25), R2 종료 \textbf{09-25}(+10) & 0 → \textbf{−10} & 규제 Must 전진 배치, 420 FP 이월 \\ \hline
A07 규제 검증 & 165 → 180 (FNLT 09-11) & 08-24~09-11 완료, 142/142 & 0 → 0 & \textbf{제약 충족} — R2 항목 재배치의 대가 \\ \hline
A08 3PL API·폴백 & 120 → 160 (~08-14) & 폴백 확정 07-08, 폴백 시험 08-21 완료, 준실시간 API는 명세·구현 완료 보관 & 20 → 12 & R-03 발동 — 비임계 흡수 \\ \hline
A09 전환 매핑 & 60 → 210 (~10-23) & 진도 62\%(9/30), CR-11로 140일 → 종료 전망 10-16 & 5 → 15 & 데이터 경로 여유 확대 \\ \hline
\textbf{A10 구현 R3 S19–S24} & \textbf{180 → 240 (09-14 ~ 12-04)} & \textbf{착수 09-28(작업일 190), 전망 종료 12-18(250)} & \textbf{0 → −10} & \textbf{임계경로 −10영업일} \\ \hline
A11 정본 판정 → G2 & 210 → 230 (FNLT 11-27) & 전망 10-19~11-13 & 5 → 10 & G2 정시 전망 \\ \hline
A12 구현 R4 잔여 & 240 → 260 & CR-08 편입(+5), 이월 300 FP(+5) → 30일 & 10 → \textbf{0} & 제2 임계경로(회복 계획) \\ \hline
A13 리허설 1~3 & 230 → 275 & 전망 11-16 착수 & 5 → 10 & \\ \hline
A14 통합시험 & 240 → 270 & as-is 250 → 280 / to-be 250 → 275(SS 겹침) & 0 → −10 / \textbf{−5} & 회복 계획 1 \\ \hline
A15 P2 인터페이스 시험 & 289 → 294 (SNET 02-12) & \textbf{G1 의결로 2차 컷오버(03-27)로 분리} — MG3 선행 조건에서 제외, MC 03-05 → 시험 03-08~12 & — & R-01 대응 \\ \hline
A16 UAT·리허설·롤백 & 270 → 294 & as-is 280 → 304 / to-be 275 → 294(UAT 1차 병행) & 0 → −10 / \textbf{0} & 회복 계획 2 \\ \hline
MG3 G3 & 294 (FNLT 02-19) & \textbf{as-is 2027-03-05 / to-be 2027-02-19} & 0 → −10 / 0 & \\ \hline
A17 컷오버 & 299 (SNET 02-26) & as-is 1차 창 이탈 → 2차 창 03-12~14 / to-be 02-26 & & \\ \hline
A20 G4 & 340 → 365 (05-28) & as-is 06-11(+10) / to-be 05-28 & 0 → −10 / 0 & 허용범위 +15 내이나 \textbf{창 0 위반} \\ \hline
\fi
\end{longtable}}
\B{}{\small 작업일 0 = 2026-01-05. 기준선 BL-01 CPM 계산표 대비.\par}

\vspace{1mm}
\noindent\begin{minipage}[t]{0.58\linewidth}\vspace{0pt}
\noindent\textbf{\small 마일스톤 전망}\par\vspace{0.5mm}
{\footnotesize
\begin{tabularx}{\linewidth}[t]{|L{34mm}|L{22mm}|Y|Y|L{16mm}|L{20mm}|}\hline
\hd{마일스톤} & \hd{기준선} & \hd{as-is 전망 (회복 없음)} & \hd{to-be (회복 계획)} & \hd{허용범위} & \hd{판정} \\ \hline
\ifwbblank
\rd{7mm}\g{[G2]} & & & & & \\ \hline
\rd{7mm}\g{[리허설 3회차]} & & & & & \\ \hline
\rd{7mm}\g{[G3]} & & & & \g{[0]} & \\ \hline
\rd{7mm}\g{[컷오버 1차]} & & & & \g{[창 0]} & \\ \hline
\rd{7mm}\g{[컷오버 2차]} & & & & & \\ \hline
\rd{7mm}\g{[G4]} & & & & \g{[+15]} & \\ \hline
\rd{7mm}\g{[ES 추세 종료]} & — & & — & & \\ \hline
\else
G2 데이터 정본 & 2026-11-27 & 11-27 (A11 여유 10) & 11-27 & +5 보고 & 정시 \\ \hline
리허설 3회차·전자금융 신청 & 2027-01-29 & 01-22 & 01-22 & +5 보고 & 정시 \\ \hline
G3 & 2027-02-19 & \textbf{2027-03-05 (−10)} & 2027-02-19 (여유 0) & 0 & \textbf{예외 보고} \\ \hline
컷오버 1차(매장·주문) & 2027-02-26~28 & \textbf{이탈 → 2차 창 03-12~14} & 02-26~28 & 0 & \textbf{예외 보고} \\ \hline
컷오버 2차(물류, R-01 분리) & 2027-03-27 주말 (G1 의결) & 03-27 & 03-27 & 프로그램 & 정시 \\ \hline
G4 & 2027-05-28 & 06-11 (+10) & 05-28 & +15 (06-18) & 내 \\ \hline
ES 추세 종료 & — & 2027-07-12 (+31) & — & +15 & 초과 (추세) \\ \hline
\fi
\end{tabularx}}
\end{minipage}\hfill
\begin{minipage}[t]{0.4\linewidth}\vspace{0pt}
\noindent\textbf{\small 버퍼 소진}\par\vspace{0.5mm}
{\footnotesize
\begin{tabularx}{\linewidth}[t]{|L{22mm}|L{14mm}|Y|L{22mm}|L{22mm}|}\hline
\hd{버퍼} & \hd{크기} & \hd{소진 (데이터 일자)} & \hd{소유자} & \hd{상태} \\ \hline
\ifwbblank
\rd{9mm}\g{[합류 버퍼]} & & & \g{[프로그램]} & \\ \hline
\rd{9mm}\g{[규제 버퍼]} & & & & \\ \hline
\rd{9mm}\g{[경로 여유]} & & & & \\ \hline
\rd{9mm}\g{[활동 여유]} & & & & \\ \hline
\else
G3 합류 버퍼 & 15영업일 & as-is 시 15 전부(2차 창 03-12) / to-be 0 & 프로그램 PMO장 & \textbf{배정 요청 대상} — P1 단독 소비 불가 \\ \hline
규제 버퍼 & 22영업일 & 0 (A07 09-11 충족) & P4 & 종료 \\ \hline
데이터 경로 & 5 & −5 (CR-11로 +10 확대) & P1 PM & 여유 15 \\ \hline
A12 여유 & 10 & 10 (CR-08 5 + 이월 5) & P1 PM & 0 — 제2 임계경로 \\ \hline
\fi
\end{tabularx}}
\end{minipage}\par

\needspace{70mm}\vspace{2mm}\noindent\textbf{\small 회복 계획 — 예외 보고 \B{EX-[nn] ([제출일], 결정 시한 [ ])}{EX-01 (2026-10-02 제출, 인지 09-28 + 4영업일; 결정 시한 10-16)}}\quad\g{대안 3개 — 각각의 일정 효과·원가·리스크·결정자. 권고를 적는다.}\par\vspace{0.5mm}
{\footnotesize
\begin{tabularx}{\linewidth}{|L{26mm}|Y|L{40mm}|L{36mm}|L{40mm}|L{30mm}|}\hline
\hd{대안} & \hd{내용} & \hd{일정 효과} & \hd{원가} & \hd{리스크} & \hd{결정자} \\ \hline
\ifwbblank
\rd{14mm}1. \g{[압축·재배치]} & & & & & \\ \hline
\rd{14mm}2. \g{[버퍼 수용 / 창 이동]} & & & & & \\ \hline
\rd{14mm}3. \g{[범위 이연]} & & & & & \\ \hline
\else
\textbf{1. 압축·재배치 (권고)} & ① A14 통합시험 준비(환경·회귀 자동화·TC 확정)를 S24와 SS 겹침 → A14 30 → 25 ② UAT 1차(POS·포털)를 통합시험 후반과 병행(SS+20) → A16 24 → 19 상당(종료 02-18 유지) ③ R2 이월 420 FP를 R3 120(속도 640 전제)·A12 300으로 재배치, CR-08 A12 편입 & MG3 02-19 회복, G3 앞 여유 0, A12 TF 0(제2 임계) & +0.60(회귀 자동화 도구·시험 인력 선투입 — 컨틴전시 R-19 배정) & R-19 통합시험 결함 유출 P 40 × 1.5 = 0.60; UAT 병행으로 심각도 2 미해결 조건(G3) 압박 & 정미란(원가·품질 트레이드오프), 이재현(속도 640 약속·변경 대가 없음 확인 09-30) \\ \hline
2. 컷오버 2차 창 수용 & 매장·주문 컷오버를 2027-03-12~14로, G3 03-05 & 합류 버퍼 15 전부 소진, 물류 2차 컷오버(03-27)와 2주 간격 → 두 번 컷오버 리스크 R-14 상승 & +1.6(2차 창 재리허설·매장 통보 재작업 [추가 설정], 관리예비비 대상 — 원인이 프로그램 버퍼 배정) & R-14 P 50 → 70; 3월 성수기 잠금(설 이후) 없음 확인 필요 & \textbf{프로그램 PMO장(버퍼 배정), 정미란} \\ \hline
3. R3 범위 이연 & 옴니 기본 시나리오 중 “매장 간 재고 이동 요청”(StRS-052 일부, 약 260 FP)을 R5로 — 헌장 §13의 이연 가능 범위 확장 & A10 −10 흡수 → MG3 02-19, 대안 1의 압축 불요 & 0 (FP 기준선 유지, 인도 시점 변경) & B-05(배달수수료 회피 37억) 실현 M+24 → M+27 지연 위험; 오정근 동의 필요 & \textbf{오정근(Senior User)}, CCB(릴리스 이연, RACI 27행) \\ \hline
\fi
\end{tabularx}}
\B{\small\g{권고와 재기준선 여부: \hrulefill}\par}{\small 권고: \textbf{대안 1을 주 계획으로, 대안 3을 조건부 예비}(11-13 S22 종료 시점 속도가 620 미만이면 자동 발동)로 승인 요청. 대안 2는 프로그램 PMO장의 버퍼 배정 판단 자료로 첨부. Day2 ⑧ 항목 10 “재기준선은 CCB + 이사회, 최대 2회” — \textbf{이번은 재기준선이 아니다}(기준선 날짜 유지, 지연은 지연으로 기록).\par}

% ------------------------------------------------------------------ 9 (가로)
\clearpage
\secbar{9. 조기 신호와 흐름 지표}
\guide{지수가 나빠지기 전에 볼 수 있었던 신호 — 시점·지표·보았어야 한 것·그때 가능했던 조치. 사후 합리화가 아니라 \textbf{다음 보고 주기의 관측 항목 목록}이다. 흐름 지표(throughput·WIP·work item age·cycle time)는 EVM과 별개가 아니다 — throughput 하락은 SI EV 하락의 선행 지표다.}
{\footnotesize
\begin{tabularx}{\linewidth}{|C{6mm}|L{34mm}|L{52mm}|Y|L{70mm}|}\hline
\hd{\#} & \hd{시점} & \hd{지표} & \hd{보았어야 한 것} & \hd{가능했던 조치} \\ \hline
\ifwbblank
\rd{10mm}1 & & & & \\ \hline
\rd{10mm}2 & & & & \\ \hline
\rd{10mm}3 & & & & \\ \hline
\rd{10mm}4 & & & & \\ \hline
\rd{10mm}5 & & & & \\ \hline
\else
1 & 2026-04 성과보고 제4호 & SPI 0.784 & 열별 분해 — 상용SW −6.0(조달) vs SI 실질 0. “개발은 정시, 조달만 지연”이라고 적었어야 함 & 라이선스 이연이 A05 개발환경에 영향 없음을 확인하고 R-04 이슈 전환(IS-01) — 실제 4월 보고는 “SPI 0.78 주의”로 끝남 \\ \hline
2 & 2026-05 2주차 & R-05 트리거(참석률 70\% 미만 2주 연속) 발동, work item age 최장 12영업일(POS 화면 정의서) & 대체 결정권자 지정(등록부 R-05 조치)이 미실행 & 오정근 부재 시 매장운영팀장 전결 지정 — 실제 지정은 06-03, 3주 늦음 \\ \hline
3 & 2026-05 성과보고 제5호 & CV +2.70 급등(4월 −0.16) & “원가 절감”이 아니라 인식 기준 변경(03-13 계약 변경)의 흔적 — 조정 CV 병기 & 5월부터 발생주의 조정 행을 넣었어야 함 — 실제 8월 감리 준비 중 발견 \\ \hline
4 & 2026-08 & throughput 620·620(정상)이나 결함 발견 주당 40 → 12(⑭), 시험 인력 규제 검증 이동 & 8월 인수 FP는 규제 항목(단순 검증 항목 다수)이라 FP 대비 시험 깊이 감소 & S18 연장(3주) 결정을 8월 중순에 내렸다면 A10 지연 −10이 아니라 −5 \\ \hline
5 & 2026-09-04 & S18 3주 연장 결정 & 그 순간 A10 착수 09-28 = TF −10이 계산 가능 & 예외 보고를 09-11에 낼 수 있었음 — 실제 10-02(허용범위 “예상 시점 5영업일” 규정에는 부합) \\ \hline
\fi
\end{tabularx}}
\vspace{2mm}\noindent\textbf{\small 스프린트 흐름 지표}\par\vspace{0.5mm}
{\footnotesize\setlength{\tabcolsep}{3pt}
\begin{tabularx}{\linewidth}{|L{58mm}|C{14mm}|C{14mm}|C{14mm}|C{14mm}|C{14mm}|C{14mm}|C{14mm}|C{14mm}|C{14mm}|C{14mm}|C{14mm}|Y|}\hline
\hd{스프린트} & \hd{S7} & \hd{S8} & \hd{S9} & \hd{S10} & \hd{S11} & \hd{S12} & \hd{S13} & \hd{S14} & \hd{S15} & \hd{S16} & \hd{S17} & \hd{S18} \\ \hline
\ifwbblank
\rd{7mm}throughput (인수 FP) & & & & & & & & & & & & \\ \hline
\rd{7mm}WIP (진행 중 StRS) & & & & & & & & & & & & \\ \hline
\rd{7mm}work item age 최장 (영업일) & & & & & & & & & & & & \\ \hline
\rd{7mm}cycle time 중앙값 (영업일) & & & & & & & & & & & & \\ \hline
\rd{7mm}결함 발견 (⑭, 스프린트 기간 주차 합) & & & & & & & & & & & & \\ \hline
\else
throughput (인수 FP) & 620 & 620 & 600 & \textbf{480} & \textbf{470} & 560 & 560 & 600 & 620 & 620 & 640 & 630 \\ \hline
WIP (진행 중 StRS) & 22 & 24 & 26 & \textbf{41} & 38 & 31 & 28 & 27 & 25 & 24 & 26 & 29 \\ \hline
work item age 최장 (영업일) & 6 & 7 & 9 & \textbf{12} & 10 & 8 & 7 & 7 & 6 & 8 & 9 & 11 \\ \hline
cycle time 중앙값 (StRS 승인 → 인수, 영업일) & 14 & 15 & 16 & 21 & 22 & 18 & 17 & 16 & 15 & 15 & 16 & 17 \\ \hline
결함 발견 (⑭, 스프린트 기간 내 주차 합) & 36 & 69 & 77 & 120 & 98 & 98 & 83 & 90 & 50 & 20 & 71 & 118 \\ \hline
\fi
\end{tabularx}}
\B{\small\g{DORA 4지표(참고): 배포 빈도 [ ] / 변경 리드타임 [ ] / 변경 실패율 [ ] / 복구 시간 [ ] — 목표와 연결된 회복 계획 항목 [ ]}\par}{\small DORA 4지표(개발 → 시험 환경, 2026-09 [추가 설정]): 배포 빈도 주 2회 / 변경 리드타임(커밋 → 시험 환경) 3.1일 / 변경 실패율 12\% / 복구 시간 4시간. 통합시험(12월) 전까지 변경 실패율 10\% 이하가 목표 — 회복 계획 ①(회귀 자동화)의 효과 지표.\par}
\landscapeoff

% ------------------------------------------------------------------ 10 (세로 — A4 한 장)
\secbar{10. 성과보고서 1장 — Highlight Report \B{제[n]호}{제9호 (2026-09-30)}}
\guide{A4 한 장. RAG의 정의 — G = 허용범위 내·전망도 내 / A = 허용범위 내·전망 초과 가능(경고) / R = 초과 예상·예외 보고 제출. 모든 축이 G이면 보고할 이유가 없다. 수치 없는 RAG는 RAG가 아니다. \textbf{결정 요청이 없는 성과보고는 보고가 아니라 기록이다.}}
{\footnotesize\setlength{\tabcolsep}{2.5pt}
\begin{tabularx}{\linewidth}{|L{15mm}|C{8mm}|L{50mm}|L{30mm}|Y|}\hline
\hd{축} & \hd{RAG} & \hd{값 (데이터 일자)} & \hd{허용범위} & \hd{전망 · 근거} \\ \hline
\ifwbblank
\rd{8mm}일정 & & & & \\ \hline
\rd{8mm}원가 & & & & \\ \hline
\rd{8mm}범위 & & & & \\ \hline
\rd{8mm}품질 & & & & \\ \hline
\rd{8mm}리스크 & & & & \\ \hline
\rd{8mm}편익 & & & & \\ \hline
\rd{8mm}지속가능성 & & & & \\ \hline
\else
일정 & \textbf{R} & 임계경로 A10 −10영업일, SPI(t) 0.920, ES −0.72개월 & 컷오버 창 0 / G4 +15 & as-is G3 03-05·창 이탈 → \textbf{EX-01 제출(10-02)}. 회복 계획 승인 시 02-19 회복(여유 0) \\ \hline
원가 & A & EV 73.43 / AC 71.13, CPI 1.032(조정 0.968), \textbf{EAC4 150.86} & 집행 한도 156.00 (EAC 기준) & 여유 5.14. EAC3′ 155.04는 경고 대역. 컨틴전시 잔여 8.79(집행 1.41 = 13.8\%, 진척 50\%) \\ \hline
범위 & G & FP 누적 변동 +185(+1.5\%), 순변동 건수 +17 & ±372 / ±35건 & CCB 10-15 4건 상정, BL-02 10-23 \\ \hline
품질 & A & 스프린트 시험 결함밀도 0.132건/FP(⑭), 심각도 1: 0 / 2: 7 미해결 & 작업패키지 0.13 / 프로젝트 0.4 & 0.13 초과 → 1단계 에스컬레이션(이재현 → P1 PM, 09-28). 감리 1차 8건 중 2건 귀속 이의 \\ \hline
리스크 & A & 잔여 위협 EMV 7.26(⑬) & 8억 / 단일 2억 & 여유 0.74. R-10 0.76 상승 주시(G2). 관리예비비 대상 2.20 \\ \hline
편익 & G & B-08 12h 설계 검증 완료(TC-FRN-040~048 통과), G2 편익 전망 118억 유지 & −5\%p & R-02 잔여 112점 동의가 B-02·B-04 가맹분 전망의 변수 \\ \hline
지속가능성 & A & 보안진단 1회차 High 2건 조치 중(기한 10-16), 개인정보 검증 142/142 & High 미조치 0(G3) & 기한 내 조치 시 G \\ \hline
\fi
\end{tabularx}}
\vspace{1.5mm}
\noindent\begin{minipage}[t]{0.42\linewidth}\vspace{0pt}
{\footnotesize
\begin{tabularx}{\linewidth}[t]{|L{34mm}|Y|}\hline
\hd{핵심 수치} & \hd{값} \\ \hline
\ifwbblank
\rd{6mm}PV / EV / AC & \\ \hline
\rd{6mm}SV / CV (실질) & \\ \hline
\rd{6mm}SPI / SPI(t) / CPI / CPI(조정) & \\ \hline
\rd{6mm}EAC / VAC / 한도 여유 & \\ \hline
\rd{6mm}ES / IEAC(t) & \\ \hline
\rd{6mm}인수 FP / 결함 미해결 & \\ \hline
\rd{6mm}컨틴전시 집행·잔여 / 버퍼 & \\ \hline
\else
PV / EV / AC & 77.83 / 73.43 / 71.13 \\ \hline
SV / CV & −4.40 (실질 −2.44) / +2.30 (실질 −2.40) \\ \hline
SPI / SPI(t) / CPI / CPI(조정) & 0.943 / 0.920 / 1.032 / 0.968 \\ \hline
EAC4 / VAC / 집행 한도 여유 & 150.86 / −5.06 / 5.14 \\ \hline
ES / IEAC(t) & 8.28개월 / 18.48개월 (추세 2027-07-12) \\ \hline
인수 FP / 결함 미해결 & 7,020 (56.6\%) / 116 (심각도 2: 7) \\ \hline
컨틴전시 집행·잔여 / 합류 버퍼 & 1.41 · 8.79 / 15 (배정 요청 예정) \\ \hline
\fi
\end{tabularx}}
\end{minipage}\hfill
\begin{minipage}[t]{0.56\linewidth}\vspace{0pt}
\begin{fieldboxu}\footnotesize\setlength{\parskip}{0.4em}
\ifwbblank
\g{이번 기간 완료:}\lines{2}{6mm}\g{다음 기간 계획:}\lines{2}{6mm}\g{예외 (제출된 예외 보고·판정):}\lines{1}{6mm}\g{이슈·리스크 상위 5:}\lines{2}{6mm}
\else
\textbf{이번 기간(9월) 완료} — 규제 ① 검증 142/142(09-11), R2 종료·인수 630 FP(09-25), 3PL 폴백 시험 통과(08-21 보고), 중간감리 1차 실시(09-07~18, 지적 8건 09-25), 보안진단 1회차(High 2·Medium 9), CR-08~11 접수·영향 분석, 변경협의체 합의 3건(09-30).\par
\textbf{다음 기간(10월) 계획} — S19·S20 인수(R3 1,240 FP + 이월 120), CCB 10-15(CR-08~11, CR-07 사후), BL-02 발행 10-23, 매핑 종료 10-16(A09), 정본 판정 착수 10-19(A11), 감리 지적 대응(기한 10-16·10-30), 보안 High 2건 조치 10-16, EX-01 결정 10-16, 협의회 회신 10-23.\par
\textbf{예외} — EX-01(일정: 컷오버 창 이탈 예상) 제출 10-02, 대안 3개, 권고 대안 1 + 조건부 대안 3. 원가 예외 없음(경고).\par
\textbf{이슈·리스크 상위 5} — IS-02 R-05 대기(재발 방지 대체 결정권자), IS-07 감리 귀속 이의 2건, IS-08 보안 High 2건, R-10 마스터 변경 요청 2건 접수(G2 앞), R-19 R4 압축.
\fi
\end{fieldboxu}
\end{minipage}\par
\vspace{1.5mm}
\noindent\textbf{\small 결정 요청} \g{(누가 · 언제까지 · 무엇을)}\par\vspace{0.5mm}
{\footnotesize
\begin{tabularx}{\linewidth}{|C{6mm}|L{40mm}|L{30mm}|Y|}\hline
\hd{\#} & \hd{누가} & \hd{언제까지} & \hd{무엇을} \\ \hline
\ifwbblank
\rd{8mm}1 & & & \\ \hline
\rd{8mm}2 & & & \\ \hline
\rd{8mm}3 & & & \\ \hline
\else
1 & 정미란 스폰서 · 프로그램 PMO장 & 2026-10-16 (EX-01 10영업일) & 회복 계획 대안 1 승인(컨틴전시 R-19 0.60 배정 포함), 대안 3 조건부 발동 규칙 승인, 대안 2 대비 합류 버퍼 배정 원칙 회신 \\ \hline
2 & CCB (정미란 위원장) & 2026-10-15 & CR-08·09·10·11 의결, CR-07 사후 확인, 8월 컨틴전시 보고 누락 소급 확인 \\ \hline
3 & 정하늘 수석감리원 · 정미란 & 2026-10-08 & 감리 지적 3·4번 귀속 이의(⑭)에 대한 회신 — 발주사 미결정 원인으로 재분류 \\ \hline
\fi
\end{tabularx}}
\end{document}
